\subsection{Estado del Arte y de la Técnica}

Este apartado analiza las soluciones tecnológicas existentes, el estricto marco regulatorio chileno y el nivel de madurez de las herramientas de Inteligencia Artificial Clínica. El análisis de estas dimensiones justifica las decisiones arquitectónicas y los límites de interoperabilidad del prototipo.

\subsubsection{Prácticas clínicas y tecnologías no autorizadas}

El traspaso de pacientes (\textit{handoff}) es un momento de alto riesgo en la gestión clínica. Para estandarizar esta comunicación, los hospitales fomentan protocolos como Situación, Antecedentes, Evaluación y Recomendación (SBAR). Sin embargo, el registro de esta actividad depende frecuentemente de libros foliados o planillas impresas. La fricción de los sistemas oficiales puede empujar al personal a adoptar tecnologías no autorizadas por la institución, incluido el uso de aplicaciones de mensajería en dispositivos personales, lo que dificulta la confidencialidad y la auditoría del dato.

\subsubsection{Ecosistema de fichas clínicas y barreras de interoperabilidad}

El sector público chileno ha avanzado en la digitalización de sus procesos mediante distintos proveedores. Esto ha producido un escenario heterogéneo en el que conviven plataformas de atención primaria y sistemas de información hospitalaria (SIH) de distinta procedencia. La interoperabilidad entre estas soluciones depende de acuerdos, estándares y controles institucionales. Como el proyecto aún no cuenta con una contraparte hospitalaria, Carer-Record incorpora al prototipo una integración simulada con datos sintéticos mediante contratos FHIR y mensajes HL7v2. Esto permite evaluar la capacidad funcional sin atribuirle una conexión institucional real y mantiene la alineación con los lineamientos del Ministerio de Salud (MINSAL)~\cite{minsal_estandares_perfiles}.

\subsubsection{Marco Normativo y Privacidad de Datos}

El diseño de software de salud en Chile está sometido a un escrutinio legal que determina su arquitectura:

\begin{itemize}
    \item \textbf{Trazabilidad y derechos:} la Ley 20.584 regula los derechos y deberes asociados a la atención en salud y refuerza la necesidad de proteger el registro clínico y su acceso~\cite{ley_20584_derechos_paciente}.

    \item \textbf{Autoría:} la identificación del usuario y la marca temporal de cada evento permiten atribuir su creación y validación. Esta trazabilidad no equivale, por sí sola, a una firma electrónica certificada.

    \item \textbf{Protección de datos:} la Ley 19.628 y la Ley 21.719 establecen obligaciones para el tratamiento de datos personales y sensibles. Este marco justifica el uso exclusivo de datos sintéticos durante la fase académica y los controles de privacidad incorporados en la arquitectura~\cite{ley_19628_vida_privada,ley_21719_datos_personales}.
\end{itemize}

\subsubsection{Tecnologías de IA y procesamiento de lenguaje natural clínico}

La viabilidad de Carer-Record depende de técnicas de procesamiento de lenguaje natural (PLN). Para la transcripción se utiliza reconocimiento automático del habla (ASR), mientras que la clasificación y síntesis requieren modelos ajustados al lenguaje clínico. Debido al ruido del entorno hospitalario y al riesgo de que los modelos generativos produzcan información incorrecta, el flujo HITL es una salvaguarda central. Lyu y Pergola muestran las posibilidades y límites de los sistemas de simplificación especializados~\cite{lyu2024}; Rust et al. evidencian que los resúmenes generados aún requieren revisión experta~\cite{rust2025}; y Lazaros et al. sistematizan el valor de mantener intervención humana en sistemas de IA~\cite{lazaros2026}.
